\documentclass[12pt]{article}

%---------------------------------------------------------------
% Idioma: español si el paquete babel-spanish está instalado;
% en caso contrario, inglés para que el documento siga compilando.
%---------------------------------------------------------------
\usepackage[utf8]{inputenc}
\usepackage[T1]{fontenc}
\IfFileExists{spanish.ldf}{%
  \usepackage[spanish]{babel}%
}{%
  \usepackage[english]{babel}%
  \PackageWarning{GOV_AKADI}{babel-spanish no instalado; se usa english}%
}
\usepackage{geometry}
\usepackage{seqsplit}
\usepackage{enumitem}
\usepackage{array}
\usepackage{longtable}

\geometry{a4paper, margin=2.5cm}
\setlist{itemsep=2pt, topsep=4pt}
\setlength{\LTpre}{6pt}
\setlength{\LTpost}{6pt}

\title{
\textbf{GOV\_AKADI / AGF-AKADI}\\
\large Documento Completo para Producción Real\\
\normalsize Versión 1.2.1-CANDIDATE
}

\date{Borrador para revisión humana --- 23 de noviembre de 2026}

\begin{document}

\maketitle

%===============================================================
\section{Naturaleza y estado del documento}

Este documento describe el diseño, el alcance, los requisitos, los
controles y el plan de verificación de GOV\_AKADI / AGF-AKADI para una
eventual operación en producción real. Es un \textbf{borrador de diseño}:
no acredita implementación, validación, autorización, activación ni
operación real de ninguna capacidad.

\subsection{Etiquetas}

\begin{description}[leftmargin=4.2cm, style=nextline]
    \item[\texttt{[VIGENTE]}] Disposición ya establecida en la línea base documental.
    \item[\texttt{[PROPUESTA]}] Contenido nuevo que completa lo no descrito. No está aprobado.
    \item[\texttt{[DECISIÓN PENDIENTE]}] Elección que solo puede tomar una autoridad humana competente.
    \item[\texttt{[SIN VERIFICAR]}] Afirmación no comprobada con evidencia.
\end{description}

\subsection{Alcance de la autorización}

\texttt{[VIGENTE]} Este documento no autoriza la ejecución de ningún motor
opaco o comercial, dentro o fuera del chat. Describe el diseño y las condiciones
que deben cumplirse antes de cualquier uso. Ningún motor comercial se ha
evaluado, y la modalidad sin API permanece NO\_APLICABLE hasta superar G1 a G5.

\subsection{Vocabulario de estado}

Los términos de ausencia se usan de forma única en todo el documento:

\begin{description}[leftmargin=4.2cm, style=nextline]
    \item[NO\_VERIFICABLE] No existen medios o evidencia suficientes para comprobar un dato dentro del alcance.
    \item[NO\_DISPONIBLE] El elemento requerido no existe en la modalidad o el entorno declarados.
    \item[NO\_APLICABLE] El requisito no es exigible en la modalidad o para ese motor.
\end{description}

%===============================================================
\section{Resumen}

\begin{itemize}
    \item El sistema busca responder con utilidad, trazabilidad y control a cualquier tipo de petición.
    \item Puede usar motores transformer y no transformer, privados o comerciales opacos.
    \item En la modalidad sin API, el motor se ejecuta en infraestructura propia y debe superar cinco puertas de admisión (G1 a G5).
    \item El código se genera sin ejecutarse; la ejecución requiere un entorno aislado y autorización por tarea.
    \item Hoy ninguna capacidad del sistema está implementada ni activada. Los componentes de prueba GA-COMP-001 y GA-EVI-001 demuestran el procedimiento de evidencia, pero no acreditan capacidades del sistema. Las decisiones D-01 a D-10 condicionan el avance.
\end{itemize}

%===============================================================
\section{Objetivo de producción real}

\subsection{Declaración}

\texttt{[PROPUESTA]} El sistema debe responder de forma útil, trazable y
controlada a peticiones de cualquier tema, proceso, documento, código,
consulta o pregunta, en un entorno de chat. Utiliza uno o varios motores de
lenguaje, privados o comerciales, incluidos motores opacos de arquitectura
transformer y no transformer. No depende en tiempo de ejecución de una API
externa del proveedor del motor.

\subsection{Distinción obligatoria}

\texttt{[VIGENTE]} Estar diseñado para producción real no equivale a estar
implementado, autorizado ni operando en producción real. Cada condición
conserva su propia clasificación (Sección~\ref{sec:estados}).

%===============================================================
\section{Alcance funcional}
\label{sec:alcance}

\texttt{[PROPUESTA]} La tabla define el riesgo base de cada categoría. El
contexto puede elevarlo, nunca reducirlo sin evidencia.

\begin{longtable}{|p{3.1cm}|p{4.6cm}|p{2.1cm}|p{4.2cm}|}
\hline
\textbf{Categoría} & \textbf{Ejemplos} & \textbf{Riesgo base} & \textbf{Tratamiento} \\
\hline
\endhead
Conversación general & Explicaciones, redacción informal, ideas & Bajo & Respuesta directa con control estándar. \\
\hline
Consulta de información & Hechos, datos, normas, precios, noticias & Medio & Indicar fecha de corte; verificar hechos vigentes con fuente. \\
\hline
Documentos & Redacción, resumen, revisión, traducción & Medio & Conservar el original; no alterar fuentes; marcar supuestos. \\
\hline
Procesos & Procedimientos, planes, flujos operativos & Medio--Alto & Declarar supuestos y dependencias; no sustituir la decisión humana. \\
\hline
Código y Python & Scripts, funciones, análisis de datos, depuración & Alto & Generación sin ejecución; ejecución solo en entorno aislado autorizado (Sección~\ref{sec:codigo}). \\
\hline
Asesoría sensible & Salud, legal, financiera, seguridad personal & Alto & Información general, límites explícitos y remisión a profesional. \\
\hline
Contenido prohibido & Armas, malware, explotación de menores, otros definidos por política & No atendible & Rechazo explícito y registro del evento. \\
\hline
\end{longtable}

\subsection{Límites}

\texttt{[VIGENTE]} Ninguna respuesta se presenta como verificada sin
evidencia que la sustente. El sistema distingue entre lo que sabe, lo que
supone y lo que no puede comprobar.

%===============================================================
\section{Clases de motor}
\label{sec:motores}

\subsection{Clasificación}

\begin{longtable}{|p{3.6cm}|p{5.4cm}|p{5.4cm}|}
\hline
\textbf{Clase} & \textbf{Descripción} & \textbf{Uso típico} \\
\hline
\endhead
T1: Transformer decodificador & Generación autorregresiva de propósito general. & Conversación, documentos, código. \\
\hline
T2: Transformer codificador o codificador-decodificador & Representación y transformación de texto. & Clasificación, recuperación, resumen controlado. \\
\hline
N1: Espacio de estados (SSM) & Modelo secuencial con estado recurrente. & Generación con contexto largo. \\
\hline
N2: Recurrente clásico & Redes recurrentes de tipo LSTM o GRU. & Clasificación y etiquetado. \\
\hline
N3: Híbrido o basado en reglas & Combinación de reglas, búsqueda y modelos pequeños. & Respuestas acotadas y validación. \\
\hline
\end{longtable}

\texttt{[DECISIÓN PENDIENTE]} La clase N2 de motores coincide en nombre con
el término N2 de la línea institucional (D-01). El modelo de estados de la
Sección~\ref{sec:estados} no usa esa etiqueta, por lo que la colisión es solo
de nombre. Si se aprueba D-07, la clase pasa a denominarse M2.

\subsection{Origen y opacidad}

Cada motor se clasifica también por origen:

\begin{itemize}
    \item \textbf{Privado propio:} pesos y código bajo control de la organización.
    \item \textbf{Abierto con licencia permisiva:} pesos disponibles bajo licencia de uso.
    \item \textbf{Comercial opaco:} pesos no disponibles o no inspeccionables; solo se conoce el comportamiento observable.
\end{itemize}

\texttt{[VIGENTE]} Para un motor comercial opaco no se infiere arquitectura,
tamaño, datos de entrenamiento ni mecanismos internos. Cualquier afirmación
sobre ellos se registra como NO\_VERIFICABLE.

\subsection{Evaluación de motores opacos}

\texttt{[PROPUESTA]} La evaluación de un motor opaco es exclusivamente
conductual: entradas controladas, salidas observadas, regresión entre
versiones y pruebas adversariales. Las conclusiones se limitan a lo observado.

\subsection{Registro de motores}

\texttt{[PROPUESTA]} Cada motor autorizado tiene una entrada en un registro
con los campos siguientes. Todo campo no verificado figura como NO\_VERIFICABLE.

\begin{itemize}
    \item Identificador interno y versión.
    \item Clase (T1, T2, N1, N2, N3) y origen (privado, abierto o comercial opaco).
    \item Licencia y estado de verificación jurídica.
    \item Disponibilidad de pesos (sí, no o parcial).
    \item Hash de pesos o del artefacto, si es accesible.
    \item Longitud máxima de contexto declarada por el proveedor.
    \item Fechas de incorporación y retirada, y responsable.
    \item Referencia a la última batería de pruebas.
\end{itemize}

%===============================================================
\section{Modalidad sin dependencia de API}
\label{sec:sin_api}

\subsection{Requisito}

\texttt{[PROPUESTA]} El motor se ejecuta dentro de infraestructura controlada
por la organización. Ninguna petición del usuario se envía a un servicio
externo de inferencia en tiempo de ejecución.

\subsection{Puertas de admisión}

Un motor opera en esta modalidad solo si cumple, con evidencia, las puertas siguientes:

\begin{enumerate}[label=G\arabic*.]
    \item \textbf{Licencia.} La licencia permite la ejecución local o en infraestructura propia sin restricciones incompatibles con el uso previsto. Estado inicial: SIN VERIFICAR. \texttt{[DECISIÓN PENDIENTE: revisión jurídica.]}
    \item \textbf{Disponibilidad de pesos.} Existen los artefactos ejecutables. Si el motor solo es accesible mediante API del proveedor, la modalidad es NO\_APLICABLE para ese motor.
    \item \textbf{Integridad.} Los artefactos tienen un hash registrado y comprobable.
    \item \textbf{Recursos.} Cómputo y memoria declarados y medidos. Estado inicial: NO\_\allowbreak VERIFICABLE.
    \item \textbf{Aislamiento.} El motor no tiene salida de red hacia servicios no autorizados.
\end{enumerate}

\subsection{Actualización y cambio de motor}

\texttt{[PROPUESTA]} Ningún motor se sustituye sin: (a) evaluación en la
batería de regresión; (b) registro del cambio con hash y versión; (c)
aprobación de la autoridad competente. Un cambio silencioso de motor se
clasifica como incidente.

\subsection{Degradación}

\texttt{[PROPUESTA]} Si el motor principal no está disponible, el sistema
responde con un mensaje de servicio limitado. No cambia a otro motor salvo
que la política de enrutamiento lo autorice para esa categoría.

%===============================================================
\section{Arquitectura de referencia}
\label{sec:arquitectura}

\texttt{[PROPUESTA]} Componentes lógicos. Su existencia material no está acreditada.

\begin{longtable}{|p{3.6cm}|p{10.6cm}|}
\hline
\textbf{Componente} & \textbf{Responsabilidad} \\
\hline
\endhead
Interfaz de chat & Recibe peticiones y presenta respuestas, con indicación de límites. \\
\hline
Clasificador de riesgo & Asigna categoría y riesgo a cada petición (Sección~\ref{sec:alcance}). \\
\hline
Enrutador & Selecciona motor según categoría, riesgo y registro (Sección~\ref{sec:motores}). \\
\hline
Motores & Ejecutan la inferencia en infraestructura controlada (Sección~\ref{sec:sin_api}). \\
\hline
Política y filtros & Aplican reglas de contenido, privacidad y seguridad antes y después de la inferencia. \\
\hline
Verificador & Comprueba condiciones verificables de la respuesta: formato, citas, código y cálculos. \\
\hline
Registro de evidencias & Guarda trazas, versiones, hashes y decisiones, sin contenido sensible innecesario. \\
\hline
Almacén de contexto & Gestiona la memoria de sesión según la política de retención. \\
\hline
Supervisión & Detecta anomalías, incidentes y desviaciones de métricas. \\
\hline
\end{longtable}

%===============================================================
\section{Flujo de una petición}
\label{sec:flujo}

\texttt{[PROPUESTA]} Flujo mínimo:

\begin{enumerate}
    \item Recepción y normalización de la entrada.
    \item Detección de instrucciones incrustadas en los datos (\emph{prompt injection}). Se tratan como datos, no como órdenes.
    \item Clasificación de categoría y riesgo.
    \item Aplicación de la política de contenido. Si el contenido no es atendible, se rechaza y se registra.
    \item Enrutamiento al motor autorizado.
    \item Inferencia.
    \item Verificación de condiciones comprobables.
    \item Aplicación de la política de salida y etiquetado de incertidumbre.
    \item Registro de la traza y entrega de la respuesta.
\end{enumerate}

\texttt{[VIGENTE]} Un fallo en cualquier paso produce una respuesta de
limitación explícita. Nunca se entrega una respuesta incompleta sin aviso.

%===============================================================
\section{Código y Python}
\label{sec:codigo}

\subsection{Generación}

\texttt{[PROPUESTA]} El sistema puede generar código en Python y otros
lenguajes. Cada fragmento incluye propósito, entradas y salidas,
dependencias con versión fijada y supuestos. No se presenta como verificado
si no se ha ejecutado en el entorno de prueba.

\subsection{Ejecución}

\texttt{[PROPUESTA]} El código solo se ejecuta en un entorno aislado, sin
acceso a datos de producción ni a red, con límites de tiempo y memoria. Cada
ejecución requiere autorización explícita del usuario.

\subsection{Estándares mínimos para Python}

\begin{itemize}
    \item Versión mínima declarada (por ejemplo, 3.11). \texttt{[DECISIÓN PENDIENTE: D-06.]}
    \item Estilo verificable con análisis estático: formato, tipado y linter. \texttt{[PROPUESTA]}
    \item Pruebas unitarias para funciones con lógica no trivial.
    \item Manejo explícito de errores; sin captura genérica silenciosa.
    \item Sin credenciales, claves ni rutas personales en el código.
    \item Dependencias fijadas en un archivo de bloqueo.
\end{itemize}

\subsection{Límites}

\texttt{[VIGENTE]} No se genera ni optimiza código malicioso: malware,
explotación de vulnerabilidades de terceros o evasión de controles. Esta
restricción no admite excepciones por contexto declarado.

%===============================================================
\section{Documentos y procesos}

\texttt{[PROPUESTA]}

\begin{itemize}
    \item El documento original no se modifica sin registro de versión.
    \item Los resúmenes y traducciones se marcan como derivados y referencian la fuente.
    \item Los procesos se describen con entradas, salidas, responsables y puntos de decisión humana.
    \item Las decisiones que afecten a personas, dinero o derechos se proponen para revisión humana. Nunca se ejecutan automáticamente.
\end{itemize}

%===============================================================
\section{Seguridad, privacidad y datos}

\subsection{Amenazas principales}

\begin{longtable}{|p{3.8cm}|p{5.0cm}|p{5.2cm}|}
\hline
\textbf{Amenaza} & \textbf{Descripción} & \textbf{Control propuesto} \\
\hline
\endhead
Inyección de instrucciones & Texto en los datos que intenta cambiar el comportamiento. & Separación de datos e instrucciones; filtros; pruebas adversariales. \\
\hline
Fuga de contexto & Exposición de información de otra sesión o de la configuración. & Aislamiento de sesiones; pruebas de regresión. \\
\hline
Datos sensibles & Tratamiento de datos de salud, identidad o finanzas sin base. & Minimización; no almacenar por defecto; avisos. \\
\hline
Alucinación de fuentes & Citas o datos inventados presentados como verificados. & Verificación de referencias; etiquetado de incertidumbre. \\
\hline
Deriva de motor & Cambio de comportamiento sin cambio registrado. & Hash, versión y regresión obligatoria. \\
\hline
Exfiltración por ejecución & Código que accede a recursos no autorizados. & Sandbox sin red; límites; revisión. \\
\hline
\end{longtable}

\subsection{Datos}

\texttt{[PROPUESTA]} La retención de sesión por defecto se limita a la
duración de la conversación. Cualquier retención adicional requiere base
documentada y consentimiento. Las categorías de datos sensibles se tratan
conforme a la normativa aplicable. \texttt{[DECISIÓN PENDIENTE: D-05.]}

%===============================================================
\section{Calidad y evaluación}

\subsection{Dimensiones}

\texttt{[PROPUESTA]} La calidad se evalúa en: exactitud factual, utilidad,
seguridad, cumplimiento de formato, coherencia con la categoría, robustez
frente a inyección y honestidad sobre los límites.

\subsection{Criterios}

Cada criterio tiene un identificador con la forma
\texttt{C-<área>-<número>-v<versión>} y se define antes de ejecutar la
prueba. Un criterio sin definición previa no sustenta ningún veredicto.

\subsection{Umbrales}

\texttt{[DECISIÓN PENDIENTE: D-04.]} Los umbrales numéricos no se fijan en
este documento. Los establece la autoridad competente con base en la
evaluación de riesgo de cada categoría.

\subsection{Conjuntos de prueba}

\texttt{[PROPUESTA]} Conjuntos independientes para funcionalidad, regresión,
adversarial, integridad de evidencias y comportamiento ante fallos. Los
fixtures sintéticos validan la estructura de los registros, no miden la
calidad de un motor.

%===============================================================
\section{Evidencia}
\label{sec:evidencia}

\texttt{[VIGENTE]} Toda evidencia es reproducible, fechada, atribuible,
versionada e identificada como REAL o SIMULACION. Un resultado simulado
nunca se presenta como resultado real.

\texttt{[PROPUESTA]} Campos mínimos de un registro de prueba: identificador,
criterio, modo de registro, estado de ejecución, veredicto, estado de
evidencia, referencia y hash (si existe), incidencias, responsable de
revisión y fecha de ejecución (vacía si no hubo ejecución real).

\texttt{[VIGENTE]} Un hash prueba correspondencia de bytes. No prueba autoría ni fecha.

%===============================================================
\section{Estados institucionales y transiciones}
\label{sec:estados}

\subsection{Modelo de estados}

\texttt{[VIGENTE]} Cada capacidad recorre una única secuencia ascendente:

\begin{longtable}{|c|p{3.7cm}|p{3.4cm}|p{5.1cm}|}
\hline
\textbf{Orden} & \textbf{Estado} & \textbf{Condición} & \textbf{Evidencia mínima} \\
\hline
\endhead
0 & DEFINIDO & --- & Especificación con identificador de requisito. \\
\hline
1 & DISEÑADO & Diseño arquitectónico & Documento de arquitectura versionado. \\
\hline
2 & IMPLEMENTADO & Implementación material & Referencia verificable al artefacto (repositorio, versión y hash). \\
\hline
3 & DISPONIBLE & Disponibilidad técnica & Registro fechado de accesibilidad. \\
\hline
4 & VALIDADO & Validación y verificación & Informe de verificador independiente. \\
\hline
5 & AUTORIZADO & Autorización aplicable & Acto de autoridad competente identificada. \\
\hline
6 & ACTIVADO & Activación operativa & Registro de activación con alcance y autorización referenciada. \\
\hline
7 & OBSERVADO & Operación real observada & Registro de ejecución real con resultados. \\
\hline
\end{longtable}

\subsection{Reglas de transición}

\begin{enumerate}[label=\arabic*.]
    \item Las transiciones ascendentes son secuenciales. No se admiten saltos.
    \item Cada transición requiere la evidencia del estado destino.
    \item Las transiciones descendentes (suspensión, revocación, reversión) son admisibles y se registran.
    \item Ningún estado se hereda por similitud entre capacidades.
\end{enumerate}

\subsection{Estado vigente}

\begin{longtable}{|p{6.5cm}|p{6.5cm}|}
\hline
\textbf{Elemento} & \textbf{Estado} \\
\hline
\endhead
Versión del documento & 1.2.1-CANDIDATE \\
\hline
Objetivo de producción real & Objetivo de diseño; no es un estado alcanzado \\
\hline
Cobertura documental & Declarada; no verificada \\
\hline
Evaluación de motores opacos & No realizada \\
\hline
Componentes de prueba GA-COMP-001 y GA-EVI-001 & IMPLEMENTADO como artefacto de prueba; no es capacidad del sistema \\
\hline
Validación de GA-EVI-001 & Pendiente: requiere dos validadores independientes \\
\hline
Modalidad sin API & NO\_APLICABLE hasta superar G1 a G5 \\
\hline
Implementación de capacidades del sistema & NO\_VERIFICABLE \\
\hline
Pruebas reales & No acreditadas \\
\hline
Aprobación AGF & No emitida \\
\hline
Autorización de producción & No emitida \\
\hline
Producción real activada & No \\
\hline
\end{longtable}

%===============================================================
\section{Procedimiento de evidencia y reproducibilidad}
\label{sec:procedimiento}

\texttt{[VIGENTE]} Reglas aplicadas en la práctica desde la revisión 1.2.0:

\begin{enumerate}[label=P\arabic*.]
    \item La especificación se fija, con archivo y hash, antes de escribir el componente.
    \item Los valores esperados se calculan con una herramienta independiente del componente probado, por ejemplo \texttt{sha256sum}.
    \item El ejecutor lee los casos desde la especificación y no contiene casos propios.
    \item Cada registro incluye, antes de la ejecución, los hashes de la especificación, del componente y del ejecutor.
    \item Un registro no se modifica después de calcular su hash. Un cambio posterior genera un incidente y un nuevo registro.
    \item Los bloques con comandos se guardan en un archivo o se ejecutan desde variables; no se pegan como texto en la terminal.
    \item Un resultado distinto del esperado se registra como FALLIDA para ese caso. La especificación no se modifica para que pase.
    \item Una validación requiere un validador distinto del diseñador, en un entorno distinto, y la obtención de los hashes de referencia.
\end{enumerate}

\texttt{[VIGENTE]} Las reglas P5 y P6 surgieron de un caso concreto: al
añadir una línea a un registro ya calculado, su hash dejó de coincidir. El
registro original se conservó como incidente y no se usó como evidencia.

\subsection{Componentes de prueba}

\begin{longtable}{|p{3.4cm}|p{4.6cm}|p{6cm}|}
\hline
\textbf{Componente} & \textbf{Función} & \textbf{Estado y pendiente} \\
\hline
\endhead
GA-COMP-001 & Conversión a mayúsculas; prueba del flujo de definición, ejecución y registro & IMPLEMENTADO como artefacto de prueba. No acredita capacidad alguna. \\
\hline
GA-EVI-001 & Verificación de SHA-256 frente a hash declarado; 5 casos & IMPLEMENTADO como artefacto de prueba. VALIDADO pendiente de dos validadores independientes. \\
\hline
\end{longtable}

\subsection{Referencias de GA-EVI-001}

\begin{longtable}{|p{4.6cm}|p{10.6cm}|}
\hline
\textbf{Artefacto} & \textbf{SHA-256} \\
\hline
\endhead
paquete\_GA-EVI-001.tgz (verificar antes de extraer) & {\footnotesize\ttfamily \seqsplit{fedfe34680ae60f3c4cfb38a5c1a617bb6f22cc8e4146332dbe8eb334098651a}} \\
\hline
spec\_GA-EVI-001.txt & {\footnotesize\ttfamily \seqsplit{68f25ff61b68789628fd9d1c177193b63874214f5d5a568c0fd1ee6a6cffa7d4}} \\
\hline
evi001.py & {\footnotesize\ttfamily \seqsplit{cf725d300b957db46eeade2639b550e78d4ba4c64acfe2a4b00dade57b92dfe3}} \\
\hline
evi001\_run.py & {\footnotesize\ttfamily \seqsplit{fc06fb127546d9f43e32f39541f98db3f6dd14baa7568b4219811d75c24e318a}} \\
\hline
casos/ok.txt & {\footnotesize\ttfamily \seqsplit{9ab1c0d4456fb24c875ada273cc004ec0f6e560def4211167926b721325925e6}} \\
\hline
casos/mal.txt & {\footnotesize\ttfamily \seqsplit{02089bfd7c9a6609fa54edf125fc75f0dd0c1a6262b443707f9c1396b5a7306e}} \\
\hline
casos/vacio.txt & {\footnotesize\ttfamily \seqsplit{e3b0c44298fc1c149afbf4c8996fb92427ae41e4649b934ca495991b7852b855}} \\
\hline
\end{longtable}

\texttt{[VIGENTE]} La ejecución de reproducción 1, realizada en el mismo
dispositivo y sobre la carpeta de origen, no cuenta como validación
independiente.

\subsection{Criterio de validación}

\texttt{[PROPUESTA]} GA-EVI-001 pasa a VALIDADO solo si se cumplen todas estas condiciones:

\begin{enumerate}
    \item Dos validadores obtienen el hash del paquete y los seis hashes de los artefactos extraídos.
    \item Ambos obtienen spec\_sha256 igual a 68f25ff6\ldots a7d4.
    \item Ambos obtienen cinco PASS y VEREDICTO\_CONJUNTO=SUPERADA.
    \item Ningún validador diseñó los casos ni escribió el componente.
    \item Cada validador trabajó en un entorno distinto: otro dispositivo o usuario.
\end{enumerate}

Si un validador obtiene un resultado distinto, se registra como incidente.
No se corrige la especificación ni se repite la prueba hasta entender la
diferencia.

\section{Operación}

\subsection{Monitorización}

\texttt{[PROPUESTA]} Seguimiento de tasas de rechazo, respuestas limitadas,
incidencias de integridad, latencia y cambios de motor. Los umbrales de
alerta son decisión pendiente.

\subsection{Incidentes}

\texttt{[PROPUESTA]} Cada incidente registra fecha, descripción, impacto,
mitigación, responsable y cierre. Estados de ciclo: Abierto, Contenido, En
verificación y Cerrado. El cierre requiere comprobación independiente.

\subsection{Reversión}

\texttt{[PROPUESTA]} Cada cambio de motor, política o código tiene un
procedimiento de reversión documentado y probado en un entorno de prueba.

%===============================================================
\section{Gobernanza y roles}

\texttt{[PROPUESTA]} Roles funcionales. La asignación de personas es decisión pendiente.

\begin{itemize}
    \item \textbf{Custodio:} responsable de la integridad documental y del registro de cambios.
    \item \textbf{Autoridad competente:} emite autorizaciones. \texttt{[DECISIÓN PENDIENTE: D-02.]}
    \item \textbf{Verificador independiente:} comprueba resultados sin haber diseñado ni implementado lo evaluado.
    \item \textbf{Responsable jurídico:} evalúa licencias y marco de datos. \texttt{[DECISIÓN PENDIENTE: D-03 y D-05.]}
\end{itemize}

%===============================================================
\section{Riesgos}

\begin{longtable}{|p{3.8cm}|p{1.9cm}|p{7.8cm}|}
\hline
\textbf{Riesgo} & \textbf{Severidad} & \textbf{Mitigación propuesta} \\
\hline
\endhead
Incumplimiento de licencia de motor comercial opaco & ALTA & Puerta G1 y revisión jurídica antes de cualquier uso. \\
\hline
Motor sin pesos accesibles tratado como local & ALTA & Puerta G2; estado NO\_APLICABLE. \\
\hline
Respuesta presentada como verificada sin evidencia & ALTA & Etiquetado de incertidumbre; verificación de condiciones comprobables. \\
\hline
Ejecución de código fuera de sandbox & ALTA & Aislamiento sin red y autorización por tarea. \\
\hline
Cambio silencioso de motor & MEDIA & Registro de versión y regresión obligatoria. \\
\hline
Retención excesiva de datos & MEDIA & Retención limitada por defecto. \\
\hline
Registro modificado tras calcular su hash & MEDIA & Reglas P5 y P6; registro de incidente y nuevo registro. \\
\hline
Validación no independiente presentada como validada & ALTA & Criterio de validación de la Sección~\ref{sec:procedimiento} y decisión D-09. \\
\hline
\end{longtable}

%===============================================================
\section{Decisiones pendientes}

\begin{longtable}{|p{1.3cm}|p{6.4cm}|p{5.6cm}|}
\hline
\textbf{ID} & \textbf{Decisión o brecha} & \textbf{Requisito para cerrar} \\
\hline
\endhead
D-01 & Expansión y autoridad de MUDF, GACF, N2, N3 y AGF & Definición documental autorizada. \\
\hline
D-02 & Autoridad competente para autorizaciones & Acto de designación. \\
\hline
D-03 & Licencias de los motores comerciales opacos & Revisión jurídica con el documento de licencia. \\
\hline
D-04 & Umbrales de aceptación por categoría & Decisión de autoridad competente con evaluación de riesgo. \\
\hline
D-05 & Marco jurídico de datos y plazos de retención & Revisión jurídica. \\
\hline
D-06 & Versión mínima de Python y estándares de estilo & Decisión técnica documentada. \\
\hline
D-07 & Nomenclatura de la clase N2 frente al estado N2 & Acuerdo de nomenclatura. \\
\hline
D-08 & Recursos de hardware y medición real & Medición con evidencia REAL. \\
\hline
D-09 & Selección independiente de casos para GA-EVI-001 & Revisión de la lista de casos por un validador distinto del diseñador; cambios en una especificación v2. \\
\hline
D-10 & Política ante incidentes de integridad de registros & Procedimiento aprobado que defina la conservación, la clasificación y la comunicación de incidentes. \\
\hline
\end{longtable}

%===============================================================
\section{Hoja de ruta condicionada}

\texttt{[PROPUESTA]} Las fases dependen de condiciones de entrada, no de fechas.

\begin{enumerate}
    \item \textbf{Fase A: cierre documental.} Entrada: D-01 y D-02 resueltas.
    \item \textbf{Fase B: evaluación de motores.} Entrada: fase A, puertas G1 a G5 documentadas para al menos un motor y, si la evaluación usa comprobación de integridad, GA-EVI-001 en estado VALIDADO.
    \item \textbf{Fase C: pruebas controladas.} Entrada: fase B y criterios con umbrales (D-04).
    \item \textbf{Fase D: decisión de activación.} Entrada: evidencia REAL verificada por un verificador independiente y autorización emitida.
\end{enumerate}

%===============================================================
\section{Regla de no autocertificación}

\texttt{[VIGENTE]} La existencia de este documento, su nivel de detalle, su
coherencia interna o la presencia de controles descritos no constituyen por
sí mismos aprobación AGF, incorporación contractual, implementación, pruebas
ejecutadas, verificación independiente, autorización de producción ni
activación. Ningún estado se eleva por afirmación documental.

%===============================================================
\section{Glosario}

\begin{description}[leftmargin=4.2cm, style=nextline]
    \item[Transformer] Arquitectura basada en atención. Puede ser decodificadora, codificadora o ambas.
    \item[No transformer] Cualquier arquitectura no basada en atención, como SSM, recurrente o híbrida.
    \item[Motor opaco] Motor cuyos pesos o mecanismos internos no son inspeccionables.
    \item[SIMULACION] Registro de un caso hipotético o sintético.
    \item[REAL] Registro de una interacción o ejecución que efectivamente ocurrió.
    \item[NO\_VERIFICABLE, NO\_DISPONIBLE, NO\_APLICABLE] Términos de ausencia definidos en la Sección~1.
    \item[Paquete] Archivo comprimido con los artefactos de una prueba, referenciado por su hash.
    \item[Validador] Persona distinta del diseñador que repite una prueba en un entorno distinto.
    \item[Registro] Salida de una ejecución, con los hashes de todo lo que intervino, que no se modifica después de calcularse su hash.
\end{description}

%===============================================================
\section*{Control de cambios}

\begin{longtable}{|p{3cm}|p{2.6cm}|p{8.6cm}|}
\hline
\textbf{Versión} & \textbf{Fecha} & \textbf{Cambio} \\
\hline
\endhead
1.0.0-CANDIDATE & Previa & Modelo de estados y regla de no autocertificación. \\
\hline
1.1.0-CANDIDATE & 23/11/2026 & Documento completo: alcance, clases de motor, modalidad sin API (G1 a G5), arquitectura, código, seguridad, evaluación, operación, riesgos y decisiones pendientes. \\
\hline
1.1.1-CANDIDATE & 23/11/2026 & Revisión de formato y claridad. Unificación del vocabulario de ausencia (NO\_VERIFICABLE sustituye a NO\_VERIFICADO). Renombrada la etiqueta de texto no comprobado a [SIN VERIFICAR]. Carga condicional de babel. Encabezados de tabla repetibles. Resumen ejecutivo. \\
\hline
1.2.0-CANDIDATE & 23/11/2026 & Estado real de los componentes de prueba y de su validación. Procedimiento de evidencia y reproducibilidad (P1 a P8). Criterio de validación para GA-EVI-001. Alcance de la autorización. Decisiones D-09 y D-10. Nuevos riesgos. \\
\hline
1.2.1-CANDIDATE & 23/11/2026 & Correcciones de consistencia: referencia de la clase N2; tabla de referencias con los seis artefactos extraídos y el paquete; cortes de línea en hashes; anchos de tablas; eliminación de xcolor sin uso. \\
\hline
\end{longtable}

\end{document}
